% ======================================================================
\section{Extension below 6~GHz}
\label{sec:below-6ghz}
% ======================================================================

The preceding analysis rests on the thin-skin approximation
$\alpha L \gg 1$, which permits replacing the exponential depth profile by
a surface delta function (\cref{eq:collapse}).  This approximation is
excellent above 6\,GHz, where the SAR penetration depth
$\delta_{\SAR} = 1/(2\alpha)$ is a fraction of a millimetre and all
transmitted power is absorbed within the top few per cent of the 10\,g
cube.  Below 6\,GHz, the penetration depth becomes comparable to or larger
than the cube side $L = 21.5\;\mathrm{mm}$, three new complications arise,
and the surface-only reduction fails.  This section quantifies the
transition, replaces the homogeneous half-space with a layered tissue
model, and derives the standing-wave SAR that results from internal
reflections.


% ----------------------------------------------------------------------
\subsection{Frequency transition}
\label{sec:freq-transition}
% ----------------------------------------------------------------------

\Cref{tab:skin-freq} collects the dielectric properties of skin tissue
from the Gabriel parametric model \cite{Gabriel1996} and the IT'IS tissue
database across the range 0.9--60\,GHz.  Three dimensionless quantities
govern the transition:

\begin{itemize}[nosep]
\item \textbf{Depth parameter} $2\alpha L$, comparing the SAR penetration
  depth to the cube side length.
\item \textbf{Depth efficiency} $h(2\alpha L)$, the fraction of the
  surface SAR that survives mass-averaging to depth~$L$.
\item \textbf{Capture fraction} $C = 1 - e^{-2\alpha L}$, the fraction of
  transmitted power absorbed within the cube depth.
\end{itemize}

\begin{table}[ht]
\centering
\caption{Skin tissue parameters across frequency.  Complex refractive
  index $\ntilde = n - j\kappa$, field attenuation rate
  $\alpha = k_0 \kappa$, SAR depth
  $\delta_{\SAR} = 1/(2\alpha)$, depth parameter $2\alpha L$ for the
  10\,g cube ($L = 21.5\;\mathrm{mm}$), depth efficiency $h(u)$, and
  capture fraction $C$.  Dielectric data: Gabriel~1996 /
  IT'IS.}
\label{tab:skin-freq}
\begin{tabular}{@{}r r r r r r r r r@{}}
\toprule
$f$ (GHz) & $\varepsilon_r$ & $\sigma$ (S/m) & $n$ & $\kappa$
  & $\alpha$ ($\mathrm{m^{-1}}$) & $2\alpha L$ & $h(2\alpha L)$ & $\delta_{\SAR}$ (mm) \\
\midrule
0.9  & 40.9 & 0.87  & 6.53 & 1.330 &   25  &  1.1 & 0.61 & 19.9 \\
2.45 & 38.0 & 1.46  & 6.22 & 0.860 &   44  &  1.9 & 0.45 & 11.3 \\
5    & 35.8 & 3.06  & 6.05 & 0.909 &   95  &  4.1 & 0.24 &  5.3 \\
10   & 31.3 & 8.01  & 5.73 & 1.256 &  263  & 11.3 & 0.088 &  1.9 \\
28   & 16.6 & 25.8  & 4.47 & 1.851 & 1086  & 46.7 & 0.021 &  0.46 \\
60   &  7.98 & 36.4 & 3.28 & 1.663 & 2092  & 89.9 & 0.011 &  0.24 \\
\bottomrule
\end{tabular}
\end{table}

\noindent
Three regimes emerge:

\paragraph{Thin skin ($\alpha L \gg 1$, above ${\sim}\,6\;\mathrm{GHz}$).}
The capture fraction $C \approx 1$: virtually all transmitted power is
absorbed within the cube depth.  The depth efficiency $h \ll 1$ and
$h(2\alpha L) \approx 1/(2\alpha L)$, so the mass-averaged SAR is simply
$\Sab/(\rhom L)$.  The energy conservation shortcut
(\cref{eq:SAR-energy}) applies directly.

\paragraph{Transition ($\alpha L \sim 1$, approximately
  1--6\,$\mathrm{GHz}$).}
The capture fraction drops below unity: at
$2.45\;\mathrm{GHz}$, $C = 0.85$, meaning 15\% of transmitted power
passes through the cube and is absorbed deeper in the body.  The depth
efficiency $h \approx 0.45$ is neither close to zero nor close to unity.
The mass-averaged SAR depends jointly on the surface distribution
\emph{and} the depth profile.  The thin-skin shortcut overpredicts
$\langle\SAR\rangle$ by $1/h \approx 2\times$; the half-space model
remains valid but the full depth integral must be retained.

\paragraph{Deep penetration ($\alpha L \ll 1$, below
  ${\sim}\,500\;\mathrm{MHz}$).}
The SAR varies slowly across the cube depth ($h \to 1$) and the
mass-averaged SAR approaches the surface SAR: $\langle\SAR\rangle \approx
(2\alpha/\rhom)\,\Sab$.  But now $\delta_{\SAR} \gg L$, so the
homogeneous half-space model itself breaks down: the field traverses
multiple tissue layers within the cube, and reflections from internal
interfaces (fat--muscle, bone--tissue) create standing waves that dominate
the SAR distribution.

\begin{remark}[Why the transition matters]
The transition regime 1--6\,GHz is precisely the frequency range where
ICNIRP~2020 switches from SAR-based to $\Sab$-based limits.  The
surface-only reduction of the previous sections provides
a clean analytical tool for the $\Sab$ regime.  Extending it below
6\,GHz requires accounting for finite penetration depth and internal
reflections, which is the subject of the remainder of this section.
\end{remark}


% ----------------------------------------------------------------------
\subsection{Layered tissue model}
\label{sec:layered}
% ----------------------------------------------------------------------

The homogeneous half-space treats tissue as a single medium with a single
attenuation constant.  Real tissue beneath the skin surface consists of
distinct layers with markedly different dielectric properties.  At
sub-6\,GHz frequencies the SAR penetration depth is comparable to or
larger than these layer thicknesses, so internal reflections are no longer
negligible.

We adopt a canonical three-layer planar model (normal incidence):

\begin{enumerate}[nosep]
\item \textbf{Skin} (epidermis + dermis): thickness
  $d_1 \approx 2\;\mathrm{mm}$, complex refractive index $\ntilde_1$,
  conductivity $\sigma_1$, mass density~$\rho_1$.
\item \textbf{Subcutaneous fat}: thickness
  $d_2 \approx 5$--$20\;\mathrm{mm}$ (varies by body region and
  individual), $\ntilde_2$, $\sigma_2$, $\rho_2$.
\item \textbf{Muscle} (semi-infinite): $\ntilde_3$, $\sigma_3$, $\rho_3$.
\end{enumerate}

\noindent
The dielectric contrast between these layers is extreme.  At
$2.45\;\mathrm{GHz}$:

\begin{center}
\begin{tabular}{@{}l c c c c c@{}}
\toprule
Tissue & $\varepsilon_r$ & $\sigma$ (S/m) & $n$ & $\kappa$ &
  $\delta_{\SAR}$ (mm) \\
\midrule
Skin   & 38.0 & 1.46 & 6.22 & 0.860 & 11.3 \\
Fat    &  5.3 & 0.08 & 2.31 & 0.127 & 76.5 \\
Muscle & 52.7 & 1.74 & 7.31 & 0.873 & 11.2 \\
\bottomrule
\end{tabular}
\end{center}

\noindent
Fat is nearly transparent ($\delta_{\SAR} = 76.5\;\mathrm{mm}$), while
skin and muscle are lossy.  This means the electromagnetic field can
traverse a thick fat layer with little attenuation and encounter the
fat--muscle interface, where the large impedance mismatch ($|r_{23}|
\approx 0.52$) generates a strong backward-travelling wave.

\paragraph{Generalized Fresnel reflection coefficients.}
Define the normal-incidence Fresnel reflection coefficient at the
interface between layers $j$ and $j+1$:
\begin{equation}\label{eq:rjj1}
  r_{j,\,j+1}
  = \frac{\ntilde_j - \ntilde_{j+1}}{\ntilde_j + \ntilde_{j+1}},
\end{equation}
and the one-way propagation factor across layer~$j$:
\begin{equation}\label{eq:Pj}
  P_j = e^{-jk_0\ntilde_j\,d_j}
  = e^{-\alpha_j d_j}\,e^{-j\beta_j d_j},
  \qquad
  |P_j| = e^{-\alpha_j d_j} \leq 1,
\end{equation}
where $\alpha_j = k_0\kappa_j$ and $\beta_j = k_0 n_j$.
The generalized (effective) reflection coefficient at each interface is
built recursively from the bottom up \cite{Chew1995}:

\begin{align}
  \widetilde{\Gamma}_3
    &= r_{2,3},
    \label{eq:Gamma3} \\[4pt]
  \widetilde{\Gamma}_2
    &= \frac{r_{1,2} + \widetilde{\Gamma}_3\,P_2^{\,2}}
            {1 + r_{1,2}\,\widetilde{\Gamma}_3\,P_2^{\,2}},
    \label{eq:Gamma2} \\[4pt]
  \widetilde{\Gamma}_1
    &= \frac{r_{0,1} + \widetilde{\Gamma}_2\,P_1^{\,2}}
            {1 + r_{0,1}\,\widetilde{\Gamma}_2\,P_1^{\,2}},
    \label{eq:Gamma1}
\end{align}

\noindent
where $r_{0,1} = (1 - \ntilde_1)/(1 + \ntilde_1)$ is the air-to-skin
coefficient and $P_j^{\,2} = e^{-2jk_0\ntilde_j d_j}$ is the round-trip
phase--attenuation factor across layer~$j$.  The overall power reflection
coefficient is $|\widetilde{\Gamma}_1|^2$, and the power transmission
coefficient replacing $\Tzero$ is
\begin{equation}\label{eq:T-layered}
  T_{\mathrm{lay}} = 1 - |\widetilde{\Gamma}_1|^2.
\end{equation}

\paragraph{Field in each layer.}
In layer~$j$ ($j = 1, 2, 3$), with $z$ measured as the local depth from
the top of that layer, the electric field takes the form
\begin{equation}\label{eq:Ej-layer}
  E_j(z) = A_j\bigl[e^{-jk_j z} + g_j\,e^{+jk_j z}\bigr],
\end{equation}
where $k_j = k_0\ntilde_j$ is the complex wave number in layer~$j$,
$A_j$ is the forward-wave amplitude at $z = 0$ (top of layer~$j$), and
\begin{equation}\label{eq:gj-def}
  \boxed{%
    g_j = \widetilde{\Gamma}_{j+1}\,P_j^{\,2}
    = \widetilde{\Gamma}_{j+1}\,e^{-2jk_j d_j}
  }
\end{equation}
is the \textbf{backward-to-forward amplitude ratio} at the top of
layer~$j$.  Physically, $g_j$ encodes the total reflected field returning
from all deeper interfaces (including multiple reflections), attenuated
and phase-shifted by the round trip across layer~$j$.

For muscle (layer~3, semi-infinite): $g_3 = 0$ because there is no
reflected wave returning from infinity.

\begin{remark}[Amplitude linkage between layers]
The forward amplitudes are related by the transmission coefficient and
propagation factor at each interface.  Continuity of the tangential
electric field at the top of layer~$j+1$ (bottom of layer~$j$) gives
\begin{equation}\label{eq:Aj-recursion}
  A_{j+1}(1 + g_{j+1})
  = A_j\bigl[e^{-jk_j d_j} + g_j\,e^{+jk_j d_j}\bigr]
  = A_j\,e^{-jk_j d_j}\bigl[1 + g_j\,e^{2jk_j d_j}\bigr],
\end{equation}
from which $A_{j+1}$ follows once $A_1$ is known from the incident field
and the overall transmission through the air--skin interface.
\end{remark}


% ----------------------------------------------------------------------
\subsection{Standing-wave SAR}
\label{sec:standing-wave}
% ----------------------------------------------------------------------

The field intensity in layer~$j$ follows from \cref{eq:Ej-layer}:
\begin{equation}\label{eq:E2-layer}
  |E_j(z)|^2
  = |A_j|^2\bigl|e^{-jk_j z} + g_j\,e^{+jk_j z}\bigr|^2.
\end{equation}
Expanding the squared modulus (with $k_j = \beta_j - j\alpha_j$ in the
engineering $e^{+j\omega t}$ convention), the three terms are:

\begin{align}
  \bigl|e^{-jk_j z} + g_j\,e^{+jk_j z}\bigr|^2
  &= e^{-2\alpha_j z}
   + |g_j|^2\,e^{+2\alpha_j z}
   + 2|g_j|\cos(2\beta_j z + \psi_j),
   \label{eq:E2-three-terms}
\end{align}

\noindent
where $\psi_j = \arg(g_j)$.  The SAR in layer~$j$ is therefore

\begin{equation}\label{eq:SAR-standing}
  \boxed{%
  \SAR_j(z)
  = \frac{\sigma_j\,|A_j|^2}{2\rho_j}
    \Bigl[
      \underbrace{e^{-2\alpha_j z}}_{\text{forward}}
    + \underbrace{|g_j|^2\,e^{+2\alpha_j z}}_{\text{backward}}
    + \underbrace{2|g_j|\cos(2\beta_j z + \psi_j)}_{\text{interference}}
    \Bigr].
  }
\end{equation}

\noindent
The three terms have distinct physical origins:

\begin{enumerate}[nosep]
\item \textbf{Forward decay.}  The incident wave attenuating
  exponentially into the layer, as in the homogeneous half-space model.
  This is the only term when $g_j = 0$ (no reflections), recovering
  the single-exponential SAR of \cref{eq:SAR-depth}.

\item \textbf{Backward growth.}  The wave reflected from deeper
  interfaces, growing exponentially toward the surface (it has
  already been attenuated on its inward journey and is now propagating
  outward through layer~$j$).  The factor $|g_j|^2$ is the round-trip
  power attenuation of the reflected wave.

\item \textbf{Interference.}  The coherent beating between the forward
  and backward waves, producing a standing-wave pattern with spatial
  period $\pi/\beta_j = \lambda_0/(2n_j)$.  At $2.45\;\mathrm{GHz}$ in
  skin, this period is $\pi/319.6 \approx 9.8\;\mathrm{mm}$.  The
  interference oscillation can shift the SAR maximum away from the
  surface if the backward wave is sufficiently strong.
\end{enumerate}

\begin{remark}[Recovery of the half-space limit]
When there are no internal reflections ($g_j = 0$ for all~$j$),
\cref{eq:SAR-standing} reduces to $\SAR = (\sigma|A|^2/2\rho)\,
e^{-2\alpha z}$, which is \cref{eq:SAR-depth} with the identification
$|A|^2 = (4\alpha/\sigma)\,\Sab$.  The layered model is therefore a
strict generalization of the half-space theory.
\end{remark}


% ----------------------------------------------------------------------
\subsection{Integrated SAR over the cube}
\label{sec:layered-integral}
% ----------------------------------------------------------------------

The 10\,g cube extends from $z = 0$ (the body surface) to $z = L$.  When
$L$ exceeds $d_1$ or $d_1 + d_2$, the cube spans multiple tissue layers.
The absorbed energy in each layer segment is computed by integrating
\cref{eq:SAR-standing} analytically.

\paragraph{Finite layer ($g_j \neq 0$, thickness $d_j$).}
The integral of $\SAR_j \rho_j$ over the full extent of layer~$j$ is

\begin{equation}\label{eq:layer-integral}
  \boxed{%
  \int_0^{d_j}\!\SAR_j\,\rho_j\,\mathrm{d}z
  = \frac{\sigma_j |A_j|^2}{2}
    \biggl[
      \frac{1 - e^{-2\alpha_j d_j}}{2\alpha_j}
    + |g_j|^2\,\frac{e^{2\alpha_j d_j} - 1}{2\alpha_j}
    + \frac{|g_j|}{\beta_j}\Bigl(
        \sin(2\beta_j d_j + \psi_j) - \sin\psi_j
      \Bigr)
    \biggr].
  }
\end{equation}

\noindent
The three terms correspond to the forward, backward, and interference
contributions respectively.  Each has a clear physical interpretation:

\begin{itemize}[nosep]
\item The forward term $\propto (1 - e^{-2\alpha_j d_j})/(2\alpha_j)$ is
  the energy absorbed from the incident wave.  For thick layers
  ($\alpha_j d_j \gg 1$) it saturates at $1/(2\alpha_j)$.
\item The backward term $\propto |g_j|^2(e^{2\alpha_j d_j} - 1)/(2\alpha_j)$
  is the energy deposited by the reflected wave as it propagates upward
  through the layer.  It can be large when $|g_j|$ is significant and the
  layer is thick.
\item The interference term oscillates with the phase traversal
  $2\beta_j d_j$ and can be positive or negative, adding or subtracting
  from the total depending on the standing-wave phase.
\end{itemize}

\paragraph{Semi-infinite layer ($g_3 = 0$, remaining depth $d_3$).}
For the muscle layer, which has no reflected wave, the integral from the
top of layer~3 over the remaining cube depth
$d_3 = L - d_1 - d_2$ simplifies to

\begin{equation}\label{eq:muscle-integral}
  \int_0^{d_3}\!\SAR_3\,\rho_3\,\mathrm{d}z
  = \frac{\sigma_3\,|A_3|^2}{2}\;\frac{1 - e^{-2\alpha_3 d_3}}{2\alpha_3}.
\end{equation}

\paragraph{Total mass-averaged SAR.}
The peak spatial-average SAR over the cube is the sum of the layer
contributions divided by the total tissue mass:
\begin{equation}\label{eq:psSAR-layered}
  \langle\SAR\rangle_{\mathrm{cube}}
  = \frac{1}{m}\sum_{j=1}^{3}
    \int_0^{\min(d_j,\,z_j^{\max})}\!\SAR_j\,\rho_j\,\mathrm{d}z,
\end{equation}
where $z_j^{\max}$ is the remaining cube depth available for layer~$j$
after subtracting the thicknesses of the layers above.  If the layer
densities are approximately equal ($\rho_1 \approx \rho_2 \approx \rho_3
\approx \rhom$), then $m \approx \rhom L^3$ and the formula simplifies.
In practice, fat density ($\rho_2 \approx 920\;\mathrm{kg/m^3}$) differs
from skin and muscle ($\rho_1, \rho_3 \approx 1100\;\mathrm{kg/m^3}$) by
roughly 15\%, which introduces a corresponding correction to the mass
normalization.

\begin{remark}[Comparison with the homogeneous result]
In the half-space model, the analogous expression is simply
$(2\alpha/\rhom)\,h(2\alpha L)\,\Sab$.  The layered formula
\cref{eq:psSAR-layered} replaces this single term with a sum of
three-term integrals per layer, each depending on the local tissue
properties and the backward-to-forward ratio $g_j$.  The additional
complexity is the price of accounting for internal reflections.
\end{remark}


% ----------------------------------------------------------------------
\subsection{Subsurface SAR peaks}
\label{sec:subsurface-peaks}
% ----------------------------------------------------------------------

In the homogeneous half-space model, the SAR maximum always occurs at the
body surface ($z = 0$), because the exponential $e^{-2\alpha z}$ is
monotonically decreasing.  With internal reflections this is no longer
guaranteed.  The backward-travelling wave can constructively interfere
with the forward wave at some depth inside the tissue, creating a SAR
maximum below the surface.

\paragraph{Criterion for a subsurface peak in layer~1.}
Differentiating \cref{eq:SAR-standing} in layer~1:
\begin{equation}\label{eq:dSAR-dz}
  \frac{\mathrm{d}\SAR_1}{\mathrm{d}z}\bigg|_{z=0}
  = \frac{\sigma_1 |A_1|^2}{2\rho_1}
    \bigl[
      -2\alpha_1
      + 2\alpha_1 |g_1|^2
      - 4\beta_1 |g_1|\sin\psi_1
    \bigr].
\end{equation}
For the SAR to increase with depth at the surface, the derivative must be
positive.  Dividing by $2\alpha_1 > 0$:

\begin{equation}\label{eq:subsurface-condition}
  \boxed{%
    |g_1|^2 - \frac{2\beta_1}{\alpha_1}\,|g_1|\sin\psi_1 > 1
    \qquad\Longleftrightarrow\qquad
    \text{SAR maximum is below the surface.}
  }
\end{equation}

\noindent
This condition involves two factors:

\begin{itemize}[nosep]
\item \textbf{Reflection strength} ($|g_1|^2$).  For this term alone to
  exceed unity, the backward wave at the surface must carry more power
  than the forward wave, which is uncommon.  Typical values at
  $2.45\;\mathrm{GHz}$ are $|g_1| \approx 0.6$, so $|g_1|^2 \approx 0.37
  < 1$.

\item \textbf{Interference sign} ($-\sin\psi_1$).  When
  $\psi_1 = \arg(g_1)$ is negative (which is typical, since the
  propagation phase $-2k_1 d_1$ rotates $\arg(g_1)$ into the lower half
  of the complex plane), $\sin\psi_1 < 0$ and the interference term
  $-(2\beta_1/\alpha_1)|g_1|\sin\psi_1$ becomes positive.  Because
  $\beta_1/\alpha_1 = n_1/\kappa_1 \approx 7$ for skin at
  $2.45\;\mathrm{GHz}$, even a modest reflection $|g_1| \approx 0.6$ can
  produce a large interference contribution.
\end{itemize}

\paragraph{Numerical evaluation at $2.45\;\mathrm{GHz}$.}
For a skin--fat--muscle stack with $d_1 = 2\;\mathrm{mm}$,
$d_2 = 10\;\mathrm{mm}$, the computed backward-to-forward ratio is
$|g_1| \approx 0.61$, $\psi_1 \approx -61^\circ$.  The subsurface
condition evaluates to
\[
  |g_1|^2 - \frac{2\beta_1}{\alpha_1}\,|g_1|\sin\psi_1
  \;=\; 0.37 + 7.7
  \;=\; 8.0 \;\gg\; 1.
\]
The condition is satisfied by a wide margin, confirming that the SAR
maximum at $2.45\;\mathrm{GHz}$ can lie below the skin surface.

\paragraph{The role of the fat layer.}
The fat layer plays a distinctive role in creating subsurface SAR
peaks.  Its conductivity is very low ($\sigma_{\mathrm{fat}} \approx
0.08\;\mathrm{S/m}$ at $2.45\;\mathrm{GHz}$, compared to
${\sim}\,1.5\;\mathrm{S/m}$ for skin and muscle), so the SAR within the
fat itself is small.  But the field amplitude can remain high throughout
the fat because of the low attenuation ($\alpha_{\mathrm{fat}} d_2
\approx 0.07$ for a $10\;\mathrm{mm}$ layer, corresponding to only
13\% power loss across the layer).

\begin{theorem}[Subsurface peak mechanism]\label{thm:subsurface}
At sub-GHz to low-GHz frequencies with thin skin over thick
subcutaneous fat over muscle, the peak SAR within the 10\,g cube can
exceed the surface SAR.  Two effects cooperate:
\begin{enumerate}[nosep]
\item The field partially traverses the low-loss fat layer without
  substantial attenuation.
\item The fat--muscle interface reflection ($|r_{2,3}| \approx 0.52$ at
  $2.45\;\mathrm{GHz}$) creates a backward wave that interferes
  constructively with the forward wave near the skin--fat boundary.
\end{enumerate}
The constructive interference produces SAR enhancements of order
$(1 + |g_1|)^2/(1 - |g_1|)^2$ at the standing-wave antinodes, which can
reach factors of 2--4 relative to the pure-decay prediction.
\end{theorem}

\begin{remark}[Implication for the surface-based reduction]
The existence of subsurface SAR peaks means that the surface quantity
$\Sab$ alone does not bound the volumetric SAR.  The reduction
$\psSAR = f(\Sab)$ of \cref{thm:main,thm:cube} requires the
homogeneous half-space assumption ($g = 0$, monotonic depth decay).
Below 6\,GHz, the layered model must be used, and the ``input'' to the
SAR computation is no longer $\Sab$ on the body surface but rather the
full set of forward amplitudes $\{A_j\}$ and backward-to-forward ratios
$\{g_j\}$ determined by the tissue stratigraphy.
\end{remark}


% ----------------------------------------------------------------------
\subsection{When the theory breaks}
\label{sec:limitations-low-freq}
% ----------------------------------------------------------------------

The layered half-space model of the preceding subsections extends the
analysis to the low-GHz range but still relies on a locally planar,
stratified geometry.  Three phenomena invalidate this picture at
sufficiently low frequencies.

\paragraph{Whole-body resonance (below ${\sim}\,300\;\mathrm{MHz}$).}
When the free-space wavelength becomes comparable to the body height
($\lambda/2 \approx 0.85\;\mathrm{m}$ at ${\sim}\,175\;\mathrm{MHz}$
for an adult, resonance peaked near $70\;\mathrm{MHz}$
for a grounded standing adult \cite{Durney1986}), the human body acts as
a resonant dipole antenna.  Internal fields are dominated by whole-body
current distributions, not by the surface-transmitted plane wave.
Localized hotspots form at anatomical constrictions (ankles, wrists, neck)
where current density is highest, and these hotspot locations bear no
relation to the surface absorbed power density.  The surface-to-volume
reduction is fundamentally inapplicable in this regime.

\paragraph{Finite-thickness structures.}
Thin anatomical features such as the ear pinna (cartilage thickness
${\sim}\,2\;\mathrm{mm}$) violate the semi-infinite half-space
assumption even at higher frequencies.  At $2\;\mathrm{GHz}$,
$\alpha d \approx 0.09$ for a $2\;\mathrm{mm}$ ear, so the one-way field
attenuation is only $e^{-\alpha d} \approx 0.91$: the transmitted field
partially exits the far side of the ear and illuminates the opposite
surface.  The SAR in such structures requires a finite-slab transmission
model rather than a half-space model.

\paragraph{Multi-path and deep hotspots.}
At frequencies where the penetration depth exceeds the body diameter at
narrow cross-sections (e.g., $\delta_{\SAR} \approx 20\;\mathrm{mm}$ at
$900\;\mathrm{MHz}$ versus forearm diameter ${\sim}\,60\;\mathrm{mm}$),
the field transmitted through one side of the body can reach the opposite
side.  Reflections from the far body--air interface create a second
backward wave, and the resulting interference pattern can produce deep
hotspots at locations not predictable from the surface fields alone.  This
``multi-path through the body'' effect is distinct from the layered-tissue
standing waves discussed above: it depends on the full 3D body geometry,
not just the local tissue stratigraphy.


